\clearpage
\fancyhead[L]{\small COMET.}
\hypertarget{native-472-lower}{}
\begin{center}
\rule{0.55\linewidth}{0.4pt}\\[2pt]
Astronomische Nachrichten. Band XXII. Nr. 514. Seite 165. 166. 1844 Aug. 18.\\[2pt]
\rule{0.55\linewidth}{0.4pt}
\end{center}
\begin{center}\textsc{Gauss} an \textsc{Schumacher}. Göttingen 1844 August 18.\end{center}

Ich sende Ihnen noch ein Paar bei dem stets ungünstigen Wetter nur mit Mühe dürftig erhaltene Cometenbeobachtungen:

\begin{center}
\setlength{\tabcolsep}{4pt}
\begin{tabular}{lccc}
1844. & \shortstack{Mittlere Zeit\\in Göttingen} & \shortstack{Gerade\\Aufsteigung} & \shortstack{Nordliche\\Abweichung} \\
Aug. 11 & $9^{\rm u}\ 28^{\rm m}\ 16^{\rm s}$ & $211^\circ\ 34'\ 13''6$ & ---\quad ---\quad --- \\
 & $9\ 31\ 48$ & ---\quad ---\quad --- & $21^\circ\ 48'\ 12''6$ \\
17 & $8\ 34\ 23$ & $208\ 42\ 2.3$ & ---\quad ---\quad --- \\
 & $8\ 36\ 8$ & ---\quad ---\quad --- & $17\ 28\ 22.6$ \\
\end{tabular}
\end{center}

\Needspace{6\baselineskip}
\SourceLeaf{473}
Am 11. hatten drei dem Cometen folgende Sterne zur Vergleichung gedient, wovon der zweite mit dem in \textsc{Bessels} Zone 460 auf folgende Art.
\[
\begin{array}{c|c|c|c}
8.\,9 & 4 & 14^{\rm u}\ 7^{\rm m}\ 50^{\rm s}54 & 21^\circ\ 50'\ 21''9
\end{array}
\]
\noindent angesetzten identisch sein wird, wenn man annimmt, dass die Beobachtung nicht am vierten sondern am dritten Faden gemacht, und also die Uhrzeit irriger Weise um den Betrag des Fadenintervalls vermindert, abgedruckt ist. Meine obige Rectascension gründet sich auf diese Hypothese, deren Richtigkeit ich noch nicht habe bestätigen können, und auf die Voraussetzung, dass die von \textsc{Bessel} angewandte Reductionszahl $15^{\rm s}46$, oder die unmittelbar beobachtete Antrittszeit $14^{\rm u}\ 8^{\rm m}\ 6^{\rm s}0$ gewesen ist. Würde aber dieser zweite Stern ganz ausgeschlossen, so würde die Rectascension des Cometen um $4''2$ grösser, also zu $211^\circ\ 34'\ 17''8$ angesetzt werden müssen. Dr. \textsc{Goldschmidts} Parabel stimmt noch immer recht gut mit den Beobachtungen überein.

Am 17. August war Piazzi XIII. 280 zur Vergleichung angewandt. Die Rectascension ist aber nur Resultat einmaliger Vergleichung.

\clearpage
\hypertarget{native-473-lower}{}
\begin{center}
\rule{0.55\linewidth}{0.4pt}\\[2pt]
Astronomische Nachrichten. Band XXII. Nr. 516. Seite 189..192. 1844 October 3.\\[2pt]
\rule{0.55\linewidth}{0.4pt}
\end{center}
\begin{flushright}1844 September 1.\end{flushright}

Gestern Abend habe ich nach langer Unterbrechung wieder ein Paar Vergleichungen des Cometen mit benachbarten Fixsternen machen können. Sie waren aber nur zur Bestimmung der geraden Aufsteigung tauglich

\begin{center}
1844 August 31.\quad $8^{\rm u}\ 25^{\rm m}\ 57^{\rm s}$ m. Z.\quad ger. Aufst. $203^\circ\ 45'\ 8^{\rm u}3$
\end{center}

Dr. \textsc{Goldschmidts} Parabel stimmt damit noch auf eine Secunde genau überein.

Ich bemerke noch, dass bei \textsc{Goldschmidts} Formeln zur Berechnung der Coordinaten, wie sie in Nr. 512 der A. N. abgedruckt sind, noch der Radius Vector als Factor hinzugefügt werden muss.

\begin{flushright}Göttingen, Sept. 3.\end{flushright}

Ich schicke Ihnen noch eine Cometenbeobachtung.

\begin{center}
1844 Sept. 2.\quad $8^{\rm u}\ 27^{\rm m}\ 11^{\rm s}$\quad G. A. $=203^\circ\ 11'\ 0''2$\\[2pt]
Abw. $=+7^\circ\ 6'\ 32''1$
\end{center}

Sie gründet sich auf zwei Vergleichungen mit einem Sterne 9. Grösse aus \textsc{Bessels} Zonen 83 und 160, und mit den beiden \textsc{Piazzi}schen Sternen XIII, 194 und 208. --- Bei dem zweiten Sterne (XIII, 194) hat mir die Rectascension einige Bedenklichkeit gemacht und es scheint fast, dass eine eigene Bewegung dabei Statt findet. Der Unterschied in Zeit zwischen dem zweiten und dritten Sterne findet sich

\begin{center}
\setlength{\tabcolsep}{3pt}
\begin{tabular}{lll}
in der Hist. Céleste & $2^{\rm m}\ 26^{\rm s}00$ & \\
in \textsc{Zachs} Catalog & $2\ 26.11$ & \\
in \textsc{Piazzis} Cat. 1. Ausg. & $2\ 26.04$ & jeder Stern 4 mal \\
\hphantom{in \textsc{Piazzis} Cat. }2. Ausg. & $2\ 26.31$ & jeder Stern 10 mal beob. \\
in \textsc{Bessels} Zone 160. & $2\ 27.00$ & \\
meine beiden Drchg. geben & $2\ 27.64$ & \\
\end{tabular}
\end{center}

\Needspace{6\baselineskip}
\SourceLeaf{474}
Der Unterschied scheint also zuzunehmen, obgleich die Präcession für XIII. 208 jährlich um ein paar Hunderttheile einer Bogensecunde kleiner ist als für XIII, 194. --- Ich habe für obige Ortsbestimmung des Cometen blos die \textsc{Bessel}schen Beobachtungen der drei Sterne angewandt und das Mittel aus allen Resultaten angesetzt. Schlösse ich für die Rectascensionen das Resultat aus XIII, 194 aus, so würde hervorgehen
\[
203^\circ\ 10'\ 57''2
\]
\noindent Uebrigens ist bei der Reduction die Refraction genau berücksichtigt. --- \textsc{Goldschmidts} Parabel gibt die G. A. noch gut, aber die Declination um 1 Minute zu klein.

\begin{flushright}Göttingen, Sept. 6.\end{flushright}

Auch gestern habe ich den Cometen beobachtet, aber blos Einen Durchgang erhalten können, der auch nur für die \emph{Declination} zu benutzen war. Das unter Berücksichtigung der Rectascension beobachtete Resultat war

\begin{center}
1844 Sept. 5.\quad $7^{\rm u}\ 56^{\rm m}\ 56^{\rm s}$ m. Z.\quad Abw. des Com. $+5^\circ\ 21'\ 20''2$
\end{center}

\textsc{Goldschmidts} Parabel gibt $1'\ 5''$ weniger.

\begin{flushright}Göttingen, Sept. 8.\end{flushright}

Heute schicke ich Ihnen meine Cometenbeobachtung von vorgestern:

\begin{center}
\setlength{\tabcolsep}{3pt}
\begin{tabular}{lll}
1844 Sept. 6. & $8^{\rm u}\ 1^{\rm m}\ 53^{\rm s}$ & Gerade Aufsteig. $202^\circ\ 7'\ 6''7$ \\
 & $8\ 8\ 7$ & Nördl. Abweich. $4\ 46\ 25.6$ \\
\end{tabular}
\end{center}

Die gerade Aufsteigung beruhet auf vier, die Abweichung auf zwei Vergleichungen.

Gestern Abend war der Himmel bedeckt.
\begin{flushright}C. F. \textsc{Gauss}.\end{flushright}
